\scanpage{28}
Alors il existe un unique $\xi\in M$ tel que l'on ait
\[
(2)\qquad\xi_k=c_{\unsure{i},k-i}\xi\quad\left(=\binom{k}{i}\xi\right)\quad\text{pour }\forall k\geq N
\]
(\uline{NB} $c_{k,\ell}$ désigne le coefficient binômial $\dfrac{(k+\ell)!}{k!\ell!}$)

\uline{Dém} \textcircled{1} Réduction au cas où $\exists p$ nb premier, tel que $\unsure{n}\operatorname{id}_M$ bijectif $\forall n\in\mathbf Z$ tel que $(n,p)=1$, i.e. $M$ un $\mathbf Z_p$-module ($\mathbf Z_p\overset{\mathrm{def}}=\mathbf Z_{(p\mathbf Z)}$).
Supposons que ce soit prouvé dans ce cas, et qu'on sache qu'il existe alors $k=k(p)$ tel que $\deleted{\unsure{\xi}}c_{i,k-i}\not\equiv0\ (p)$ (i.e. $c_{k,k-i}\in\mathbf Z_p^*$). Prouvons le th. en général. Les $c_{k,k-i}$ $(k\geq N)$ ne sont pas divisibles par un $=$ nb premier, donc sont premiers dans leur ensemble, donc $\exists k_1\ldots k_r\geq N$ et des entiers $m_1,\ldots m_r\in\mathbf Z$ tels que
\[
\sum_{\alpha=1}^{r}m_\alpha c_{k_\alpha,k_\alpha-i}=1.
\]
Donc les relations (2) impliquent, en multipliant par les $m_\alpha$ et sommant
\[
\xi=\sum_1^r m_\alpha\xi_{k_\alpha}
\]
d'où l'unicité de $\xi$. Prouvons qu'avec ce choix de $\xi$, on a bien (2) i.e.
\[
\xi_k=c_{k,k-i}\sum_1^r m_\alpha\xi_{k_\alpha}
\]
Pour prouver que cette relation est vraie, \deleted{\unsure{Alors}} il suffit de prouver qu'elle est vraie dans chaque $M_p=M\otimes_{\mathbf Z}\mathbf Z_p$ $(p\in\mathbf P)$. Mais c'est une conséquence immédiate du fait que le th. est vrai pour $M_p$.
\mspage{27}
